\section{Prime closures: a control regime and a failure regime}
\label{sec:exhibit-prime}

This exhibit applies the closure audit to a classic compression of discrete structure into an analytic object: the integers and the primes into the zeta function.
The aim is not to prove anything about $\zeta$.
It is to show that the same route-mismatch machinery used for derivatives separates a regime where two staged descriptions provably converge to a common object from a regime where naive staging fails.

\subsection{Two staged descriptions}
For $\Re(s)>1$ the zeta function has an additive and a multiplicative description,
\[
  \zeta(s)=\sum_{n\ge 1}n^{-s}=\prod_p\frac{1}{1-p^{-s}},
\]
see, for example, \cite{Titchmarsh1986}.
We stage both at a cutoff $N$ with the smoothing weight $w(u)=e^{-u}$:
\begin{align*}
  S_N(s)&:=\sum_{n=1}^{N}w(n/N)\,n^{-s},\\
  P_N(s)&:=\exp\Bigl(\sum_{p\le N}w(p/N)\,\bigl(-\log(1-p^{-s})\bigr)\Bigr)=\prod_{p\le N}(1-p^{-s})^{-w(p/N)}.
\end{align*}
The Euler-side weights act on the logarithms of the factors, and $w\equiv 1$ would give the sharp partial sum and partial product.
In the language of Section~\ref{sec:dictionary}, $S_N$ is an additive staging protocol and $P_N$ a multiplicative one; a stable closure requires them to become compatible as $N$ grows.

\subsection{Packaging maps}
To compare staged functions across regions we apply three simple packaging maps.
With the completion factor $C(s)=\tfrac12 s(s-1)\pi^{-s/2}\Gamma(s/2)$, set
\begin{align*}
  \mathrm{Raw}(F)&=F, \qquad \mathrm{Comp}(F)(s)=C(s)F(s),\\
  \mathrm{Sym}(F)(s)&=\tfrac12\bigl(\mathrm{Comp}(F)(s)+\mathrm{Comp}(F)(1-s)\bigr).
\end{align*}
$\mathrm{Comp}$ is a one-sided completion and $\mathrm{Sym}$ enforces the symmetry $s\leftrightarrow 1-s$ by construction.
Neither is an analytic continuation scheme.
In the convergence region we use $\mathrm{Comp}$, because $\mathrm{Sym}$ would evaluate the truncations at $1-s$ with $\Re(1-s)<0$, far outside their domain of validity; in the critical strip we use $\mathrm{Sym}$ as a prototype of self-dual packaging.
The two regimes thus differ in both the region and the packaging map: the exhibit compares two staging--packaging pairs, not a single isolated variable.

\subsection{Diagnostics}
Let $\mathrm{Pack}$ be the packaging map of a regime, $A_N=\mathrm{Pack}(S_N)$ and $B_N=\mathrm{Pack}(P_N)$, and let $K$ be a finite test set: $\sigma\in\PrimeSigmaConvList$ in the convergence region or $\sigma\in\PrimeSigmaStripList$ in the strip, with $t\in\PrimeTList$ and $s=\sigma+it$.
With cutoffs $N\in\PrimeNList$ we compute the route mismatch
\[
  \RM_2(N)=\frac{\bigl(\sum_{s\in K}|A_N(s)-B_N(s)|^2\bigr)^{1/2}}{\bigl(\sum_{s\in K}|A_N(s)|^2\bigr)^{1/2}+\varepsilon_0}
\]
with $\varepsilon_0=10^{-30}$ (and its analogue $\RM_\infty$ with maximum norms), together with two error ledgers against a reference computed by \texttt{mpmath.zeta} under the same packaging,
\[
  \mathrm{Err}_S(N)=\frac{\|A_N-\mathrm{Pack}(\zeta)\|_2}{\|\mathrm{Pack}(\zeta)\|_2+\varepsilon_0},\qquad
  \mathrm{Err}_P(N)=\frac{\|B_N-\mathrm{Pack}(\zeta)\|_2}{\|\mathrm{Pack}(\zeta)\|_2+\varepsilon_0}.
\]
All arithmetic, including the ratios $n/N$ and $p/N$ inside the smoothing weights, is carried out at $\PrimeDPS$-digit working precision in \texttt{mpmath}; the reported summaries are stored as double-precision numbers.
In the convergence region the error ledgers are accuracy tests; in the strip they measure how far naive staging is from the analytically continued function.

\subsection{The control regime is provably convergent}
In the region $\Re(s)>1$ the finite trend has a proof behind it.

\begin{proposition}[Convergence of the staged closures]\label{prop:prime-control}
Let $K$ be a compact subset of $\{\Re(s)>1\}$. Then $S_N\to\zeta$ and $P_N\to\zeta$ uniformly on $K$ as $N\to\infty$, and the same holds after applying $\mathrm{Comp}$. Consequently $\RM_2(N)$, $\mathrm{Err}_S(N)$, and $\mathrm{Err}_P(N)$ tend to zero on every finite $K$ in this region.
\end{proposition}

\begin{proof}
Let $\sigma_0>1$ be the minimum of $\Re(s)$ on $K$. Then $|n^{-s}|\le n^{-\sigma_0}$, which is summable, while the weights $w(n/N)$ lie in $[0,1]$ and tend to $1$ for each fixed $n$; dominated convergence, applied uniformly in $s\in K$, gives $S_N\to\zeta$ uniformly. For primes $|p^{-s}|\le p^{-\sigma_0}\le 2^{-\sigma_0}<1$, so the principal logarithm is given by its power series and $|{-\log(1-p^{-s})}|\le p^{-\sigma_0}/(1-2^{-\sigma_0})$, again summable. The weighted log-sums therefore converge uniformly to $\sum_p-\log(1-p^{-s})$, and continuity of $\exp$ on the bounded range gives $P_N\to\zeta$ by the Euler product. The factor $C(s)$ is analytic, hence bounded, on $K$, so multiplication by it preserves uniform convergence. The numerators of the three diagnostics therefore tend to zero, and their denominators are bounded below by $\varepsilon_0>0$.
\end{proof}

\subsection{Results}
Table~\ref{tab:prime-rm} and Figure~\ref{fig:prime-rm} report both regimes.
In the convergence region all three diagnostics decrease with $N$; the route mismatch reaches $\RM_2(800)=\PrimeRMConvAtEightHundred$.
In the strip they behave very differently: $\RM_2$ grows from $\PrimeRMStripAtFifty$ at $N=50$ to $\PrimeRMStripAtEightHundred$ at $N=800$, the Euler-side error reaches $\PrimeErrPStripAtEightHundred$, and even the additive error grows to $\PrimeErrSStripAtEightHundred$.

% Auto-generated by scripts/export_results_tex.py from notes/math_review/
\begin{table}[htbp]
\centering
\caption{Prime closure diagnostics. Left: convergence region with one-sided completion. Right: critical strip with symmetrized completion.}
\label{tab:prime-rm}
\small
\begin{tabular}{r rrr rrr}
\toprule
& \multicolumn{3}{c}{$\Re(s)\in\PrimeSigmaConvList$} & \multicolumn{3}{c}{$\Re(s)\in\PrimeSigmaStripList$} \\
\cmidrule(lr){2-4}\cmidrule(lr){5-7}
$N$ & $\mathrm{RM}_2$ & $\mathrm{Err}_S$ & $\mathrm{Err}_P$ & $\mathrm{RM}_2$ & $\mathrm{Err}_S$ & $\mathrm{Err}_P$ \\
\midrule
50 & $0.0760$ & $0.209$ & $0.147$ & $340$ & $6.59$ & $2.12\times 10^{3}$ \\
100 & $0.0698$ & $0.172$ & $0.112$ & $6.00\times 10^{3}$ & $10.4$ & $6.03\times 10^{4}$ \\
200 & $0.0667$ & $0.142$ & $0.0848$ & $6.20\times 10^{5}$ & $16.5$ & $9.85\times 10^{6}$ \\
400 & $0.0613$ & $0.118$ & $0.0654$ & $3.14\times 10^{8}$ & $27.1$ & $8.29\times 10^{9}$ \\
800 & $0.0524$ & $0.0985$ & $0.0506$ & $5.01\times 10^{12}$ & $43.7$ & $2.17\times 10^{14}$ \\
\bottomrule
\end{tabular}
\end{table}


\begin{figure}[htbp]
\centering
\includegraphics[width=\linewidth]{fig_prime.pdf}
\caption{Prime closure diagnostics. (a) For $\Re(s)>1$, with one-sided completion, the route mismatch and both errors against $\zeta$ decrease over the tested cutoffs, consistent with the convergence to zero established by Proposition~\ref{prop:prime-control}. (b) In the critical strip, with symmetrized completion, the same naive protocols drift apart by many orders of magnitude. The vertical scales differ by design; data in Table~\ref{tab:prime-rm}.}
\label{fig:prime-rm}
\end{figure}

To check that the stored summaries are not artifacts of the working precision, the full grid was recomputed independently at $\PrimeCheckDPS$ digits.
Every reported value of $\RM_2$, $\RM_\infty$, $\mathrm{Err}_S$, and $\mathrm{Err}_P$ agrees with the $\PrimeDPS$-digit computation to double precision.
This is a stability check of double-precision summaries, not an interval certificate.

\subsection{Interpretation}
The growth in the strip has an ordinary explanation.
For real $s\in(0,1)$ the smoothed additive sums grow like $N^{1-s}$, and the weighted prime sums $\sum_{p\le N}w(p/N)\,p^{-s}$ diverge; since $-\log(1-p^{-s})\ge p^{-s}$, they are lower bounds for the logarithmic Euler sums, which therefore diverge as well.
Applying a completion factor or a symmetry projection does not repair truncation formulas that are invalid in the region.
The measurement is therefore a finite record of a feasibility failure of these staging--packaging pairs on this grid.
It is not a theorem that the mismatch blows up at every point of the strip, and it carries no information about the location of the zeros of $\zeta$.

The methodological point is that the same route-mismatch audit used for derivatives separates the two regimes, and that in the regime where it reports success there is a short proof that success is correct.
A staging scheme that is valid in the strip, such as one built on the approximate functional equation, is the natural next protocol to audit with the same tools (Section~\ref{sec:discussion}).
